\section{שיעור 08 — משוואות דיפרנציאליות חלקיות מסדר ראשון ושיטת הקרקטריסטיקות}

\begin{mainbox}{מידע כללי}
\textbf{תאריך:} 25/11/2025\\
\textbf{נושא:} משוואות דיפרנציאליות חלקיות (מד"ח) מסדר ראשון, שיטת הקרקטריסטיקות, בעיית קושי.
\end{mainbox}

\subsection{מבוא: מד"ח מסדר ראשון לעומת מד"ר}

בניגוד למשוואות דיפרנציאליות רגילות (\LR{ODEs}), בהן הפתרון הכללי מכיל קבועים שרירותיים, הפתרון הכללי של משוואה דיפרנציאלית חלקית (\LR{PDE}) מסדר ראשון מכיל \textbf{פונקציה שרירותית}.

לדוגמה, נבחין בין המקרים:
\begin{itemize}
    \item \textbf{במד"ר:} המשוואה $y' = 0$ גוררת $y(x) = C$ (קבוע שרירותי אחד).
    \item \textbf{במד"ח:} המשוואה $\frac{\partial u}{\partial x} = 0$ עבור פונקציה $u(x,y)$ גוררת $u(x,y) = f(y)$. כל פונקציה של $y$ תקיים את המשוואה, כיוון שהנגזרת שלה לפי $x$ מתאפסת. מכאן שישנה שרירותיות רבה יותר ("אינסוף דרגות חופש" פונקציונליות).
\end{itemize}

\subsection{שיטת הקרקטריסטיקות (\LR{Method of Characteristics})}

נדון במשוואות לינאריות הומוגניות מסדר ראשון מהצורה:
\begin{equation}
    a(x,y)\frac{\partial u}{\partial x} + b(x,y)\frac{\partial u}{\partial y} = 0
    \label{eq:pde_homogen}
\end{equation}

\begin{definitionbox}{אינטרפרטציה גיאומטרית}
ניתן להגדיר שדה וקטורי $\vec{A}(x,y) = (a(x,y), b(x,y))$. אזי המשוואה (\ref{eq:pde_homogen}) ניתנת לכתיבה כ:
\[ \vec{A} \cdot \nabla u = 0 \]
משמעות הדבר היא שהגרדיאנט של $u$ מאונך לוקטור $\vec{A}$ בכל נקודה. מאחר והגרדיאנט מאונך לקווי הגובה (\LR{level sets}) של הפונקציה, נובע כי \textbf{קווי הגובה של $u$ מקבילים לקווי השדה של $\vec{A}$}.
במילים אחרות: הפונקציה $u(x,y)$ נשארת \textbf{קבועה} לאורך קווי השדה של $\vec{A}$. קווים אלו נקראים \textbf{עקומות קרקטריסטיות} (\LR{Characteristic Curves}).
\end{definitionbox}

\subsubsection{מציאת הקרקטריסטיקות}
השיפוע של קווי השדה של $\vec{A}$ מקיים $\frac{dy}{dx} = \frac{b}{a}$. לכן, המשוואה המגדירה את הקרקטריסטיקות היא:
\begin{formulabox}{משוואת הקרקטריסטיקות}
\[ \frac{dx}{a(x,y)} = \frac{dy}{b(x,y)} \]
\end{formulabox}
פתרון משוואה זו מניב משפחת עקומות מהצורה $\xi(x,y) = C$. אלו הם המסלולים במרחב $xy$ שלאורכם הפתרון $u$ קבוע.

\begin{theorembox}{הפתרון הכללי}
אם המשוואה $\xi(x,y) = C$ מתארת את הקרקטריסטיקות של המשוואה $a u_x + b u_y = 0$, אזי הפתרון הכללי הוא:
\[ u(x,y) = f(\xi(x,y)) \]
כאשר $f$ היא פונקציה שרירותית גזירה.
\end{theorembox}

\subsubsection{דוגמאות לפתרון כללי}

\begin{examplebox}{דוגמה 1: הסעה לינארית}
המשוואה: $\frac{\partial u}{\partial x} - \frac{\partial u}{\partial y} = 0$.
\\
כאן $\vec{A} = (1, -1)$. משוואת הקרקטריסטיקות:
\[ \frac{dx}{1} = \frac{dy}{-1} \implies dy = -dx \implies y = -x + C \implies x + y = C \]
הקרקטריסטיקות הן קווים ישרים $x+y=C$.
\\
\textbf{הפתרון הכללי:} $u(x,y) = f(x+y)$.
\end{examplebox}

\begin{examplebox}{דוגמה 2: משוואה עם מקדמים משתנים}
המשוואה: $x \frac{\partial u}{\partial x} + y \frac{\partial u}{\partial y} = 0$.
\\
כאן $\vec{A} = (x, y)$. משוואת הקרקטריסטיקות:
\[ \frac{dx}{x} = \frac{dy}{y} \implies \ln|x| = \ln|y| + \text{const} \implies \frac{y}{x} = C \]
הקרקטריסטיקות הן קווים העוברים דרך הראשית (קרניים).
\\
\textbf{הפתרון הכללי:} $u(x,y) = f\left(\frac{y}{x}\right)$.
\end{examplebox}

\begin{examplebox}{דוגמה 3: רוטציה}
המשוואה: $y \frac{\partial u}{\partial x} - x \frac{\partial u}{\partial y} = 0$.
\\
כאן $\vec{A} = (y, -x)$. משוואת הקרקטריסטיקות:
\[ \frac{dx}{y} = \frac{dy}{-x} \implies -x dx = y dy \implies x dx + y dy = 0 \implies \frac{1}{2}x^2 + \frac{1}{2}y^2 = \text{const} \]
הקרקטריסטיקות הן מעגלים $x^2 + y^2 = C$.
\\
\textbf{הפתרון הכללי:} $u(x,y) = f(x^2 + y^2)$.
\end{examplebox}


\subsection{בעיית קושי (\LR{Cauchy Problem})}

כדי למצוא פתרון \textbf{פרטי} (כלומר, לקבוע את הפונקציה $f$), עלינו לקבל מידע נוסף. במד"ח, תנאי זה נקרא "תנאי קושי" או תנאי שפה. התנאי ניתן בדרך כלל על גבי עקומה במרחב $(x,y)$.

\begin{notebox}{הערה חשובה לגבי תנאי שפה}
לא מספיק לתת את ערך הפונקציה בנקודה בודדת (כמו במד"ר), אלא יש לתת את ערך הפונקציה לאורך עקומה החותכת את הקרקטריסטיקות. עקומה זו נקראת \textbf{עקומת ההתחלה} או \textbf{עקומת המידע}.
\end{notebox}

\begin{physicsapplication}{יישום: מציאת פתרון פרטי}
נתבונן בדוגמה 1: $u(x,y) = f(x+y)$.
\\
\textbf{תנאי קושי:} נתון כי על הישר $x=0$, הפונקציה מקיימת $u(0,y) = y^2$.
\\
\textbf{פתרון:}
נציב את התנאי בפתרון הכללי:
\[ u(0,y) = f(0+y) = f(y) \]
מצד שני נתון $u(0,y) = y^2$. מכאן נסיק כי צורת הפונקציה היא $f(t) = t^2$.
כעת נציב חזרה לפתרון הכללי עם הארגומנט המקורי:
\[ u(x,y) = (x+y)^2 \]
בדיקה: נגזרות חלקיות נותנות $2(x+y) - 2(x+y) = 0$.
\end{physicsapplication}

\begin{physicsapplication}{דוגמה נוספת}
עבור המשוואה $x u_x + y u_y = 0$ (דוגמה 2), עם פתרון כללי $u = f(y/x)$.
\\
\textbf{תנאי קושי:} על הפרבולה $y=x^2$, נתון $u(x,x^2) = e^x$.
\\
\textbf{פתרון:}
נציב בפתרון הכללי:
\[ u(x,x^2) = f\left(\frac{x^2}{x}\right) = f(x) \]
נתון שזה שווה ל-$e^x$, לכן הפונקציה $f$ היא האקספוננט $f(t) = e^t$.
הפתרון הפרטי הוא:
\[ u(x,y) = e^{y/x} \]
\end{physicsapplication}

\subsection{הרחבה לממדים גבוהים (3D)}

עבור משוואה בשלושה משתנים $u(x,y,z)$:
\[ a u_x + b u_y + c u_z = 0 \]
אנו מחפשים משטחים עליהם $u$ קבועה. משוואות הקרקטריסטיקות הן:
\begin{formulabox}{מערכת הקרקטריסטיקות ב-3D}
\[ \frac{dx}{a} = \frac{dy}{b} = \frac{dz}{c} \]
\end{formulabox}
פתרון מערכת זו דורש מציאת שני אינטגרלים ראשוניים בלתי תלויים (שני משטחים):
\[ u_1(x,y,z) = C_1, \quad u_2(x,y,z) = C_2 \]
הקרקטריסטיקות הן עקומות החיתוך של שני המשטחים הללו. הפתרון הכללי הוא פונקציה שרירותית של שני הקבועים הללו:
\[ u(x,y,z) = F(u_1, u_2) \]

\begin{examplebox}{דוגמה תלת-ממדית}
המשוואה: $x u_x + y u_y + z u_z = 0$.
\\
המשוואות הקרקטריסטיות:
\[ \frac{dx}{x} = \frac{dy}{y} = \frac{dz}{z} \]
נפתור בזוגות:
1. $\frac{dx}{x} = \frac{dy}{y} \implies \ln x = \ln y + c \implies \frac{y}{x} = C_1$.
2. $\frac{dx}{x} = \frac{dz}{z} \implies \frac{z}{x} = C_2$.
\\
הפתרון הכללי הוא פונקציה של שני המשתנים החדשים:
\[ u(x,y,z) = F\left(\frac{y}{x}, \frac{z}{x}\right) \]
לדוגמה, פונקציה ספציפית שבחר המרצה להדגמה: $F(\eta, \zeta) = \eta e^{\zeta}$. הצבה תיתן $u = \frac{y}{x}e^{z/x}$, וניתן לוודא שזהו פתרון.
\end{examplebox}

\subsection{הוכחת נכונות הפתרון הכללי (באמצעות החלפת משתנים)}

נרצה להוכיח פורמלית ש-$u(x,y) = f(\xi(x,y))$ הוא אכן הפתרון הכללי, כאשר $\xi(x,y)=C$ הן הקרקטריסטיקות.
נגדיר מערכת קואורדינטות חדשה $(z, \eta)$ (בהרצאה סומן לעתים כ-$z$ ולעתים כמשתנה אחר, נשתמש בסימון עקבי להוכחה):
נגדיר משתנה אחד להיות הקרקטריסטיקה עצמה: $z = \xi(x,y)$.
נחשב את הנגזרת של $u$ לפי המשתנים המקוריים באמצעות כלל השרשרת:
\[ \frac{\partial u}{\partial x} = \frac{du}{dz}\frac{\partial z}{\partial x}, \quad \frac{\partial u}{\partial y} = \frac{du}{dz}\frac{\partial z}{\partial y} \]
נציב במשוואה המקורית $a u_x + b u_y = 0$:
\[ a \frac{du}{dz}\frac{\partial z}{\partial x} + b \frac{du}{dz}\frac{\partial z}{\partial y} = \frac{du}{dz} \left( a \frac{\partial \xi}{\partial x} + b \frac{\partial \xi}{\partial y} \right) \]
הביטוי בסוגריים מתאפס מכיוון ש-$\xi$ קבועה לאורך הקרקטריסטיקות (נגזרת כיוונית לאורך השדה מתאפסת), ולכן המשוואה מתקיימת זהותית. בהרצאה הוצגה הוכחה דומה המראה ש-$u$ תלויה רק בקרקטריסטיקות ולא במשתנה הניצב להן.

\subsection{סיכום שלבים לפתרון}
\begin{steps}
    \item זהה את המקדמים $a(x,y)$ ו-$b(x,y)$ במשוואה $au_x + bu_y = 0$.
    \item כתוב את משוואת הקרקטריסטיקות $\frac{dx}{a} = \frac{dy}{b}$.
    \item פתור את המד"ר ומצא את הקבוע האינטגרטיבי: $\xi(x,y) = C$.
    \item כתוב את הפתרון הכללי: $u(x,y) = f(\xi(x,y))$.
    \item אם נתון תנאי קושי (שפה), הצב אותו כדי למצוא את הפונקציה $f$.
\end{steps}